\documentclass{article}
\usepackage{graphicx}
\usepackage{amsmath}
\usepackage{amsfonts}
\usepackage{bookmark}
\usepackage{booktabs}
\usepackage{exercises}
\graphicspath{ {} }
\usepackage{indentfirst}
\usepackage{amsthm}
\usepackage{yhmath}
\usepackage{geometry}
\usepackage{placeins}
\usepackage{csquotes}
\usepackage{pgfplots}
\newcommand*{\oldneg}{\mathord{\sim}}
\newcommand\showdiv[1]{\overline{\smash{)}#1}}
\newcommand\Mydiv[1]{%
\smash{\raise.3ex\hbox{$\big)$}}\mkern-8mu
        \overline{\enspace\strut#1}}
\setlength{\tabcolsep}{5pt}
\theoremstyle{definition}
\newtheorem{definition}{Definition}

\setlength{\parindent}{.5in}
\setlength{\parskip}{12pt}
\renewcommand{\arraystretch}{2}
\geometry{margin=1in}
\newcommand{\?}{\stackrel{?}{=}}
\pgfplotsset{width=7cm,compat=1.18}\usepgfplotslibrary{statistics}

\title{Chapter Review (pg. 424)}
\begin{document}
\maketitle

\section*{1}
\[ (t+3)^6=_6C_0t^6+_6C_1t^5\times 3+_6C_2t^4\times 3^2+_6C_3t^3\times 3^3+_6C_4t^2\times 3^4+_6C_5t\times 3^5+_6C_6\times 3^6 \]
\[ =t^6+18t^5+135t^4+540t^3+1215t^2+1458t+729 \]

\section*{2}
\[ (2a+b^2)^4=_4C_0(2a)^4+_4C_1(2a)^3\times b^2+_4C_2(2a)^2\times (b^2)^2+_4C_3(2a)\times (b^2)^3+_4C_4\times (b^2)^4 \]
\[ =16a^4+32a^3b^2+24a^2b^4+8ab^6+b^8 \]

\section*{3}
\[ (w-8v)^4=_4C_0w^4+_4C_1w^3\times (-8v)+_4C_2w^2\times (-8v)^2+_4C_3w\times (-8v)^3+_4C_4\times (-8v)^4 \]
\[ =w^4-32w^3v+384w^2v^2-2048wv^3+4096v^4 \]

\section*{4}
\[ (r^3-4s)^5=_5C_0(r^3)^5+_5C_1(r^3)^4\times (-4s)+_5C_2(r^3)^3\times (-4s)^2+_5C_3(r^3)^2\times (-4s)^3+_5C_4(r^3)\times (-4s)^4+_5C_5(-4s)^5 \]
\[ =r^{15}-20r^{12}s+160r^9s^2-650r^6s^3+1280r^3s^4-1024s^5 \]

\section*{5}
To find the number of samples in which there are at least 4 flavors, we need to subtract the number of combinations of samples less than 4 flavors from the total number of samples. For each flavor, because you can 
either choose it or not, there are $2^{15}$ combinations in total. Thus the total number of combinations of at least 4 samples is \[ 2^{15}-_{15}C_0+_{15}C_1+_{15}C_2+_{15}C_3=32192 \]

\section*{6} 
\[ _8C_60.5^6\times 0.5^2\approx 0.109 \]

\section*{7}
\[ _8C_40.5^4\times 0.5^4\approx 0.273 \]

\section*{8}
\[ _8C_70.5^7\times 0.5\approx 0.031 \]

\section*{9}
\[ _8C_0\times 0.5^8\approx 0.004 \]

\section*{10}
The $z$ score is $\frac{89-95}{7}\approx -0.9$ and thus $P(x\le 89)\approx .1841$.

\section*{11}
The $z$ score is $\frac{84-95}{7}\approx -1.6$ and thus $P(x\le 84)\approx .0548$.

\section*{12}
The $z$ scores for $x=91$ and $x=100$ are $\frac{91-95}{7}\approx -0.6$ and $\frac{100-95}{7}\approx 0.7$, respectively, meaning $P(x\le 91)\approx .2743$ and $P(x\le 100)\approx .7580$. Thus 
$P(91<x\le 100)\approx .7580-.2743=.4837$.

\section*{13}
The $z$ score is $\frac{50-95}{7}\approx -6.4$ and thus $P(x\le 50)=.0000+$, or slightly more than zero.

\section*{14}
The $z$ score for $x=100$ is $\frac{100-95}{7}\approx 0.7$ meaning $P(x\le 100)\approx .7580$. Thus $(P(x>100)\approx 1-.7580=.242$.

\section*{15}
We have already seen (in problem 13) that $P(x\le 50)=.0000+$. The $z$ score for $x=80$ is $\frac{80-95}{7}\approx -2.1$ meaning $P(x\le 80)\approx .0179$, and thus 
$P(50<x\le 80)\approx .0179-.0000+\approx .0179$. 

\section*{16}
\[ \pm\frac{1}{\sqrt{300}}\approx \pm 0.058 \]  or about $\pm 5.8\%$.

\section*{17}
\[ \pm\frac{1}{\sqrt{2500}}=\pm 0.020 \] or $\pm 2.0\%$.

\section*{18}
\[ \pm\frac{1}{\sqrt{800}}\approx \pm 0.035 \] or about $\pm 3.5\%$. 

\section*{19}
\[ \pm\frac{1}{\sqrt{4900}}\approx \pm 0.014 \] or about $\pm 1.4\%$.

\section*{20}
There were 517 people surveyed. Thus the margin of error was \[ \pm\frac{1}{\sqrt{517}}\approx \pm 0.044 \]  or about $\pm 4.4\%$.

\section*{21}
This is a survey.

\section*{22}
This is an observational study.

\section*{23}
The explanation given does not necessarily match any of the described scenarios.

\section*{24}
This is an observational study.

\end{document}
